%!TEX root = ED_Arvore.tex

\section{Heap}
\begin{frame}[label=notas]
	Estrutura de dados abstrata, derivada da árvore, que satisfaz a propriedade:
	\begin{itemize}
		\item Se $B$ é filho de $A$, então $B.chave \le A.chave$ (heap de máximo)\pause
		\item Se $B$ é filho de $A$, então $B.chave \ge A.chave$ (heap de mínimo)
		\item Pode ser construída em tempo linear.
	\end{itemize}
	\pause
	\begin{center}
	\begin{tikzpicture}[->]
		\Tree[ .100 [ .19 [ .17 [ .2 ] [.7 ] ] [ .3 ] ] [ .36 [.25 ] [ .1 ] ] ]
	\end{tikzpicture}
	\end{center}
	\pause
	Não há restrições quanto ao número de filhos por nó.
	\subitem{Na prática: 2}
\end{frame}

\subsection{Aplicações}
\begin{frame}[label=notas]
	\begin{itemize}[<+->]
		\item \emph{\alert<1>{É a estrutura de dados mais eficiente para implementar uma fila de prioridade.}}
		\item Heapsort: algoritmo de ordenação \emph{in-place}, complexidade $O(n{\cdot}log(n))$.
			\vspace{-1\baselineskip}
			\begin{itemize}
				\item Vantagem: O comportamento é logarítmico, qualquer que seja a entrada.
				\item Desvantagem: não é estável; o caso médio é pior que o do quicksort e a implementação mais complexa.
				\item Recomendado para aplicações que não podem tolerar um caso desfavorável.
			\end{itemize}
		\item Algoritmos de seleção (custo linear ou mesmo constante):
			\begin{itemize}
				\item Busca de máximo/mínimo.
				\item Busca do valor médio
				\item Busca do $k$-ésimo máximo elemento
			\end{itemize}
		\item Algoritmos de grafos (ex: algoritmo de Dijkstra)
	\end{itemize}
\end{frame}

\tikzstyle{vetor} = [rectangle, draw, fill]
\begin{frame}[label=notas]

	Operações comuns:
	\begin{itemize}
		\item busca-max: encontra o máximo item (ou busca-min)
		\item remove-max: remove a raiz
		\item insere: insere um novo valor
		\item fusão: une duas heaps (como uma heap)
	\end{itemize}
	\pause
	\begin{columns}[T]
		\column{.7\textwidth}
			Implementação mais simples/eficiente ([semi-]completa): \alert<2>{vetor}
		\column{.25\textwidth}
			\vspace{-2\baselineskip}
			\Tree [.$i$ [ .$2i+1$ ] [ .$2i+2$ ] ]
			\\\vfill\hfill ${\forall}i = 0,1,...,n/2$
	\end{columns}
	\pause
	\begin{center}
	\begin{tikzpicture}
		\node[vetor] (100) {\textcolor{white}{100}};
		\node[vetor, greenUnB!35, right of=100]	(19)	{\textcolor{black}{19}};
		\node[vetor, blueUnB!35, right of=19]  	(36)	{\textcolor{black}{36}};
		\node[vetor, greenUnB!70, right of=36] 	(17)	{\textcolor{white}{17}};
		\node[vetor, greenUnB!70, right of=17] 	(3) 	{\textcolor{white}{3}};
		\node[vetor, blueUnB!70, right of=3]   	(25)	{\textcolor{white}{25}};
		\node[vetor, blueUnB!70, right of=25]  	(1) 	{\textcolor{white}{1}};
		\node[vetor, greenUnB, right of=1]     	(2) 	{\textcolor{white}{2}};
		\node[vetor, greenUnB, right of=2]     	(7) 	{\textcolor{white}{7}};

		\node[below of=100] {\footnotesize 0};
		\node[below of=19]	{\footnotesize 1};
		\node[below of=36]	{\footnotesize 2};
		\node[below of=17]	{\footnotesize 3};
		\node[below of=3] 	{\footnotesize 4};
		\node[below of=25]	{\footnotesize 5};
		\node[below of=1] 	{\footnotesize 6};
		\node[below of=2] 	{\footnotesize 7};
		\node[below of=7] 	{\footnotesize 8};

		\uncover<4>{
		\draw[->] (100) .. controls +(north:.5) and +(north:.5) .. (19);
		\draw[->] (100) .. controls +(north:1) and +(north:1) .. (36);}
		\uncover<5>{
		\draw[->] (19) .. controls +(south:.5) and +(south:.5) .. (17);
		\draw[->] (19) .. controls +(south:1) and +(south:1) .. (3);
		\draw[->] (36) .. controls +(north:.5) and +(north:.5) .. (25);
		\draw[->] (36) .. controls +(north:1) and +(north:1) .. (1);}
		\uncover<6>{
		\draw[->] (17) .. controls +(south:.5) and +(south:.5) .. (2);
		\draw[->] (17) .. controls +(south:1) and +(south:1) .. (7);}
	\end{tikzpicture}
	\end{center}
\end{frame}